% Part: first-order-logic
% Chapter: sequent-calculus
% Section: proving-things

\documentclass[../../../include/open-logic-section]{subfiles}

\begin{document}

\iftag{FOL}
      {\olfileid{fol}{seq}{pro}}
      {\olfileid{pl}{seq}{pro}}

\olsection{نمونه‌هایی از \usetoken{P}{derivation}}

\begin{ex}
یک $\Log{LK}$-!!{derivation} برای دنبالهٔ
$!A \land !B \Sequent !A$ به دست دهید.

با نوشتن دنبالهٔ پایانیِ مطلوب در پایین !!{derivation} آغاز می‌کنیم.
\begin{prooftree}
\AxiomC{}
\UnaryInf$!A\land !B \fCenter !A$
\end{prooftree}
سپس باید دریابیم چه نوع استنتاجی می‌تواند دنبالهٔ پایینی‌ای به این
صورت داشته باشد. ممکن است قاعده‌ای ساختاری باشد، اما بهتر است جست‌وجو
را با یک قاعدهٔ منطقی آغاز کنیم. تنها رابط منطقیِ دنبالهٔ پایینی
$\land$ است؛ پس در پی قاعده‌ای برای $\land$ هستیم، و چون نماد
$\land$ در مقدم رخ می‌دهد، قاعدهٔ \LeftR{\land} را بررسی می‌کنیم.
\begin{prooftree}
\AxiomC{}
\RightLabel{\LeftR{\land}}
\UnaryInf$!A\land !B \fCenter !A$
\end{prooftree}
برای دنبالهٔ بالاییِ استنتاج \LeftR{\land} دو امکان وجود دارد:
دنبالهٔ بالایی می‌تواند $!A \Sequent !A$ یا $!B \Sequent !A$ باشد.
روشن است که $!A \Sequent !A$ دنباله‌ای آغازین است (و این مطلوب است)،
درحالی‌که $!B \Sequent !A$ در حالت کلی اشتقاق‌پذیر نیست. دنبالهٔ
بالایی را درج می‌کنیم:
\begin{prooftree}
\Axiom$!A \fCenter !A$
\RightLabel{\LeftR{\land}}
\UnaryInf$!A\land !B \fCenter !A$
\end{prooftree}
اکنون $\Log{LK}$-!!{derivation} درستی برای دنبالهٔ
$!A \land !B \Sequent !A$ داریم.
\end{ex}

\begin{ex}
یک $\Log{LK}$-!!{derivation} برای دنبالهٔ
$\lnot !A \lor !B \Sequent !A \lif !B$ به دست دهید.

با نوشتن دنبالهٔ پایانیِ مطلوب در پایین !!{derivation} آغاز کنید.
\begin{prooftree}
\AxiomC{}
\UnaryInf$\lnot !A \lor !B \fCenter !A \lif !B$
\end{prooftree}
برای یافتن قاعده‌ای منطقی که بتواند این دنبالهٔ پایانی را به دست دهد،
به رابط‌های منطقی آن نگاه می‌کنیم: $\lnot$، $\lor$ و $\lif$. اکنون
فقط $\lor$ و $\lif$ برای ما مهم‌اند، زیرا در دنبالهٔ پایانی نقش
!!{main operator}s را برای !!{sentence}s دارند؛ اما $\lnot$ در دامنهٔ
رابط دیگری است و بعداً به آن می‌پردازیم. پس گزینه‌های ما برای قاعدهٔ
منطقیِ استنتاج نهایی، قاعدهٔ \LeftR{\lor} و قاعدهٔ \RightR{\lif}
هستند. در واقع هر کدام را می‌توان برگزید، اما قاعدهٔ
\RightR{\lif} را برمی‌گزینیم (حتی اگر تنها دلیل این باشد که انشعاب
به دو شاخه را به تعویق می‌اندازد). بنا بر صورت استنتاج‌های
\RightR{\lif} که می‌توانند دنبالهٔ پایینی را به دست دهند، باید چنین
داشته باشیم:
\begin{prooftree}
\AxiomC{}
\UnaryInf$ !A, \lnot !A \lor !B \fCenter !B $
\RightLabel{\RightR{\lif}} \UnaryInf$ \lnot !A \lor !B \fCenter !A \lif !B $
\end{prooftree}

اگر $\lnot !A \lor !B$ را به بیرونی‌ترین جای مقدم منتقل کنیم، می‌توانیم
قاعدهٔ \LeftR{\lor} را اعمال کنیم. بنا بر شِما، این کار باید به دو
دنبالهٔ بالاییِ زیر منشعب شود:
\begin{prooftree}
\AxiomC{}
\UnaryInf$\lnot !A, !A \fCenter !B$
\AxiomC{}
\UnaryInf$!B, !A \fCenter !B$
\RightLabel{\LeftR{\lor}}
\BinaryInf$ \lnot !A \lor !B, !A \fCenter !B $
\RightLabel{\LeftR{\Exchange}}
\UnaryInf$ !A, \lnot !A \lor !B \fCenter !B $
\RightLabel{\RightR{\lif}}
\UnaryInf$ \lnot !A \lor !B \fCenter !A \lif !B $
\end{prooftree}
به یاد داشته باشید که می‌کوشیم مسیر خود را رو به بالا تا دنباله‌های
آغازین بیابیم؛ ظاهراً بسیار نزدیک شده‌ایم!{} شاخهٔ راست فقط به یک
تضعیف و یک تبادل نیاز دارد تا به دنباله‌ای آغازین برسد و کامل شود:
\begin{prooftree}
\AxiomC{}
\UnaryInf$\lnot !A, !A \fCenter !B$
\Axiom$!B \fCenter !B$
\RightLabel{\LeftR{\Weakening}}
\UnaryInf$!A, !B \fCenter !B$
\RightLabel{\LeftR{\Exchange}}
\UnaryInf$!B, !A \fCenter !B$
\RightLabel{\LeftR{\lor}}
\BinaryInf$\lnot !A \lor !B, !A \fCenter !B $
\RightLabel{\LeftR{\Exchange}}
\UnaryInf$ !A, \lnot !A \lor !B \fCenter !B $
\RightLabel{\RightR{\lif}}
\UnaryInf$ \lnot !A \lor !B \fCenter !A \lif !B $
\end{prooftree}

اکنون به شاخهٔ چپ نگاه می‌کنیم. تنها رابط منطقی در هر !!{sentence}،
نماد $\lnot$ در !!{sentence}s مقدم است؛ پس با نمونه‌ای از قاعدهٔ
\LeftR{\lnot} سروکار داریم.
\begin{prooftree}
\AxiomC{}
\UnaryInf$ !A \fCenter !B, !A$
\RightLabel{\LeftR{\lnot}}
\UnaryInf$\lnot !A, !A \fCenter !B$
\Axiom$!B \fCenter !B$
\RightLabel{\LeftR{\Weakening}}
\UnaryInf$!A, !B \fCenter !B$
\RightLabel{\LeftR{\Exchange}}
\UnaryInf$!B, !A \fCenter !B$
\RightLabel{\LeftR{\lor}}
\BinaryInf$\lnot !A \lor !B, !A \fCenter !B $
\RightLabel{\LeftR{\Exchange}}
\UnaryInf$ !A, \lnot !A \lor !B \fCenter !B $
\RightLabel{\RightR{\lif}}
\UnaryInf$ \lnot !A \lor !B \fCenter !A \lif !B $
\end{prooftree}
همانند تکمیل شاخهٔ راست، برای تکمیل این شاخهٔ چپ نیز فقط یک تضعیف
و یک تبادل فاصله داریم.
\begin{prooftree}
\Axiom$!A \fCenter !A$
\RightLabel{\RightR{\Weakening}}
\UnaryInf$ !A \fCenter !A, !B$
\RightLabel{\RightR{\Exchange}}
\UnaryInf$ !A \fCenter !B, !A$
\RightLabel{\LeftR{\lnot}}
\UnaryInf$\lnot !A, !A \fCenter !B$
\Axiom$!B \fCenter !B$
\RightLabel{\LeftR{\Weakening}}
\UnaryInf$!A, !B \fCenter !B$
\RightLabel{\LeftR{\Exchange}}
\UnaryInf$!B, !A \fCenter !B$
\RightLabel{\LeftR{\lor}}
\BinaryInf$\lnot !A \lor !B, !A \fCenter !B $
\RightLabel{\LeftR{\Exchange}}
\UnaryInf$ !A, \lnot !A \lor !B \fCenter !B $
\RightLabel{\RightR{\lif}}
\UnaryInf$ \lnot !A \lor !B \fCenter !A \lif !B $
\end{prooftree}
\end{ex}

\begin{ex}
یک $\Log{LK}$-!!{derivation} برای دنبالهٔ
$\lnot !A \lor \lnot !B \Sequent \lnot (!A \land !B)$ به دست دهید.

با استفاده از روش‌های بالا، با نوشتن دنبالهٔ پایانیِ مطلوب در پایین
آغاز می‌کنیم.
\begin{prooftree}
\AxiomC{}
\UnaryInf$ \lnot !A \lor \lnot !B \fCenter \lnot (!A \land !B) $
\end{prooftree}
رابط‌های اصلیِ در دسترسِ !!{sentence}s دنبالهٔ پایانی، نماد
$\lor$ و نماد $\lnot$ هستند. در اینجا اعمال هر یک از قاعدهٔ
\LeftR{\lor} یا قاعدهٔ \RightR{\lnot} نتیجه می‌دهد، اما با قاعدهٔ
\RightR{\lnot} آغاز می‌کنیم، زیرا انشعاب به دو شاخه را اندکی به
تعویق می‌اندازد:
\begin{prooftree}
\AxiomC{}
\UnaryInf$!A \land !B, \lnot !A \lor \lnot !B \fCenter $
\RightLabel{\RightR{\lnot}}
\UnaryInf$\lnot !A \lor \lnot !B \fCenter \lnot (!A \land !B)$
\end{prooftree}
اکنون میان بررسی قاعدهٔ \LeftR{\land} و قاعدهٔ \LeftR{\lor} حق انتخاب
داریم. ببینیم با اعمال قاعدهٔ \LeftR{\land} چه رخ می‌دهد: می‌توانیم
با دنبالهٔ $!A, \lnot !A \lor !B \Sequent \quad$ یا دنبالهٔ
$!B, \lnot !A \lor !B \Sequent \quad$ آغاز کنیم. چون !!{derivation}
نسبت به $!A$ و $!B$ متقارن است، مورد نخست را برمی‌گزینیم:
\begin{prooftree}
\AxiomC{}
\UnaryInf$!A, \lnot !A \lor \lnot !B \fCenter $
\RightLabel{\LeftR{\land}}
\UnaryInf$!A \land !B, \lnot !A \lor \lnot !B \fCenter $
\RightLabel{\RightR{\lnot}}
\UnaryInf$\lnot !A \lor \lnot !B \fCenter \lnot (!A \land !B)$
\end{prooftree}
با ادامهٔ پرکردن !!{derivation} می‌بینیم که به مشکلی برمی‌خوریم:
\begin{prooftree}
\Axiom$!A \fCenter !A$
\RightLabel{\LeftR{\lnot}}
\UnaryInf$ \lnot !A, !A \fCenter$
\AxiomC{}
\RightLabel{?}
\UnaryInf$!A \fCenter !B$
\RightLabel{\LeftR{\lnot}}
\UnaryInf$ \lnot !B, !A \fCenter$
\RightLabel{\LeftR{\lor}}
\BinaryInf$\lnot !A \lor \lnot !B, !A \fCenter $
\RightLabel{\LeftR{\Exchange}}
\UnaryInf$!A, \lnot !A \lor \lnot !B \fCenter $
\RightLabel{\LeftR{\land}}
\UnaryInf$!A \land !B, \lnot !A \lor \lnot !B \fCenter $
\RightLabel{\RightR{\lnot}}
\UnaryInf$\lnot !A \lor \lnot !B \fCenter \lnot (!A \land !B)$
\end{prooftree}
بالای شاخهٔ راست را نمی‌توان بیش از این فروکاست و با استنتاج‌های
ساختاری نیز نمی‌توان آن را به دنباله‌ای آغازین رساند؛ پس این مسیر
درستی نیست. بنابراین آشکار است که اعمال قاعدهٔ \LeftR{\land} در بالا
خطا بود. اگر به مرحلهٔ پیشین بازگردیم و به‌جای آن قاعدهٔ
\LeftR{\lor} را اعمال کنیم، به عبارت زیر می‌رسیم:
\begin{prooftree}
\AxiomC{}
\UnaryInf$\lnot !A, !A \land !B \fCenter $

\AxiomC{}
\UnaryInf$\lnot !B, !A \land !B \fCenter $

\RightLabel{\LeftR{\lor}}
\BinaryInf$\lnot !A \lor \lnot !B, !A \land !B \fCenter $
\RightLabel{\LeftR{\Exchange}}
\UnaryInf$!A \land !B, \lnot !A \lor \lnot !B \fCenter $
\RightLabel{\RightR{\lnot}}
\UnaryInf$\lnot !A \lor \lnot !B \fCenter \lnot (!A \land !B)$
\end{prooftree}
با تکمیل هر شاخه به همان شیوهٔ پیشین، داریم:
\begin{prooftree}
\Axiom$ !A \fCenter!A$
\RightLabel{\LeftR{\land}}
\UnaryInf$!A \land !B \fCenter !A$
\RightLabel{\LeftR{\lnot}}
\UnaryInf$\lnot !A, !A \land !B \fCenter $

\Axiom$ !B \fCenter !B$
\RightLabel{\LeftR{\land}}
\UnaryInf$!A \land !B \fCenter !B$
\RightLabel{\LeftR{\lnot}}
\UnaryInf$\lnot !B, !A \land !B \fCenter $

\RightLabel{\LeftR{\lor}}
\BinaryInf$\lnot !A \lor \lnot !B, !A \land !B \fCenter $
\RightLabel{\LeftR{\Exchange}}
\UnaryInf$!A \land !B, \lnot !A \lor \lnot !B \fCenter $
\RightLabel{\RightR{\lnot}}
\UnaryInf$\lnot !A \lor \lnot !B \fCenter \lnot (!A \land !B)$
\end{prooftree}
(در این گام‌ها می‌توانستیم قواعد $\land$ را پایین‌تر از قواعد
$\lnot$ اعمال کنیم و باز هم !!{derivation} درستی به دست آوریم.)
\end{ex}

\begin{ex}
تا اینجا از قاعدهٔ انقباض استفاده نکرده‌ایم، اما گاهی به آن نیاز
داریم. نمونه‌ای از چنین وضعی را ببینیم. فرض کنید می‌خواهیم
$\quad \Sequent !A \lor \lnot !A$ را اثبات کنیم. اعمال وارونهٔ
$\RightR{\lor}$ یکی از !!{derivation}s زیر را به ما می‌دهد:
\begin{prooftree}
\AxiomC{}
\UnaryInf$ \fCenter !A$
\RightLabel{\RightR{\lor}}
\UnaryInf$ \fCenter !A \lor \lnot !A$
\DisplayProof\qquad\bottomAlignProof
\AxiomC{}
\UnaryInf$!A \fCenter $
\RightLabel{\RightR{\lnot}}
\UnaryInf$ \fCenter \lnot !A$
\RightLabel{\RightR{\lor}}
\UnaryInf$ \fCenter !A \lor \lnot !A$
\end{prooftree}
البته هیچ‌یک از این دو به دنباله‌ای آغازین ختم نمی‌شود. نکته این است
که قاعدهٔ انقباض به ما امکان می‌دهد دو نسخه از !!a{sentence} را در
یکی ادغام کنیم---و هنگامی که در جست‌وجوی برهانیم، یعنی از پایین به
بالا می‌رویم، می‌توانیم نسخه‌ای از $!A \lor \lnot !A$ را در مقدمه
نگه داریم؛ برای نمونه:
\begin{prooftree}
\AxiomC{}
\UnaryInf$ \fCenter !A \lor \lnot !A, !A$
\RightLabel{\RightR{\lor}}
\UnaryInf$ \fCenter !A \lor \lnot !A, !A \lor \lnot !A$
\RightLabel{\RightR{\Contraction}}
\UnaryInf$ \fCenter !A \lor \lnot !A$
\end{prooftree}
اکنون می‌توانیم $\RightR{\lor}$ را بار دوم اعمال کنیم و
~$\lnot !A$ را نیز به دست آوریم؛ این به !!{derivation} کاملی می‌انجامد.
\begin{prooftree}
\Axiom$!A \fCenter !A$
\RightLabel{\RightR{\lnot}}
\UnaryInf$\fCenter !A, \lnot !A$
\RightLabel{\RightR{\lor}}
\UnaryInf$\fCenter !A, !A \lor \lnot !A$
\RightLabel{\RightR{\Exchange}}
\UnaryInf$ \fCenter !A \lor \lnot !A, !A$
\RightLabel{\RightR{\lor}}
\UnaryInf$ \fCenter !A \lor \lnot !A, !A \lor \lnot !A$
\RightLabel{\RightR{\Contraction}}
\UnaryInf$ \fCenter !A \lor \lnot !A$
\end{prooftree}
\end{ex}

\begin{prob}
برای دنباله‌های زیر !!{derivation}sی لازم را به دست دهید:
\begin{enumerate}
\item $!A \land (!B \land !C) \Sequent (!A \land !B) \land !C$.
\item $!A \lor (!B \lor !C) \Sequent (!A \lor !B) \lor !C$.
\item $!A \lif (!B \lif !C) \Sequent !B \lif (!A \lif !C)$.
\item $!A \Sequent \lnot\lnot !A$.
\end{enumerate}
\end{prob}

\begin{prob}
برای دنباله‌های زیر !!{derivation}sی لازم را به دست دهید:
\begin{enumerate}
\item $(!A \lor !B) \lif !C \Sequent !A \lif !C$.
\item $(!A \lif !C) \land (!B \lif !C) \Sequent (!A \lor !B) \lif !C$.
\item $\Sequent \lnot(!A \land \lnot !A)$.
\item $!B \lif !A \Sequent \lnot !A \lif \lnot !B$.
\item $\Sequent (!A \lif \lnot !A) \lif \lnot !A$.
\item $\Sequent \lnot(!A \lif !B) \lif \lnot !B$.
\item $!A \lif !C \Sequent \lnot (!A \land \lnot !C)$.
\item $!A \land \lnot !C \Sequent \lnot (!A \lif !C)$.
\item $!A \lor !B, \lnot !B \Sequent !A$.
\item $\lnot !A \lor \lnot !B \Sequent \lnot(!A \land !B)$.
\item $\Sequent (\lnot !A \land \lnot !B) \lif\lnot(!A \lor !B)$.
\item $\Sequent \lnot(!A \lor !B) \lif (\lnot !A \land \lnot !B)$.
\end{enumerate}
\end{prob}

\begin{prob}
برای دنباله‌های زیر !!{derivation}sی لازم را به دست دهید:
\begin{enumerate}
\item $\lnot(!A \lif !B) \Sequent !A$.
\item $\lnot(!A \land !B) \Sequent \lnot !A \lor \lnot !B$.
\item $!A \lif !B \Sequent \lnot !A \lor !B$.
\item $\Sequent \lnot \lnot !A \lif !A$.
\item $!A \lif !B, \lnot !A \lif !B \Sequent !B$.
\item $(!A \land !B) \lif !C \Sequent (!A \lif !C) \lor (!B \lif !C)$.
\item $(!A \lif !B) \lif !A \Sequent !A$.
\item $\Sequent (!A \lif !B) \lor (!B \lif !C)$.
\end{enumerate}
(همهٔ این موارد به قاعدهٔ~$\RightR{\Contraction}$ نیاز دارند.)
\end{prob}
\end{document}
